% !TEX program = xelatex
% Самодостаточный конспект. Компилировать XeLaTeX/LuaLaTeX (tectonic 06-bilinejnye-kvadratichnye-formy.tex).
\documentclass[11pt,a4paper]{article}
\newcommand{\coursename}{Специальные разделы высшей математики}
\newcommand{\rightheader}{§6. Билинейные и квадратичные формы}
\newcommand{\doctitle}{Билинейные и квадратичные формы}
% ==== Общая преамбула конспектов (собирается в каждый .tex как самодостаточная) ====
% Перед этим файлом должны быть определены: \documentclass и макросы
%   \coursename  — название курса (левый колонтитул, подзаголовок),
%   \rightheader — правый колонтитул (напр. «§2. Рациональные дроби»),
%   \doctitle    — заголовок раздела на титуле.
\usepackage{fontspec}
\setmainfont{cmunrm.otf}[
  BoldFont       = cmunbx.otf,
  ItalicFont     = cmunti.otf,
  BoldItalicFont = cmunbi.otf]
\usepackage{polyglossia}
\setdefaultlanguage{russian}
\setotherlanguage{english}

\usepackage{amsmath,amssymb,amsthm,mathtools}
\usepackage[a4paper,margin=2.2cm]{geometry}
\usepackage{enumitem}
\usepackage{graphicx}
\usepackage{array}
\usepackage{xcolor}
\definecolor{accent}{HTML}{1F6FEB}
\definecolor{rulecol}{HTML}{D0D7DE}
\usepackage[colorlinks=true,linkcolor=accent,urlcolor=accent,citecolor=accent]{hyperref}
\usepackage{titlesec}
\usepackage{fancyhdr}

% ----- Колонтитулы -----
\pagestyle{fancy}
\fancyhf{}
\fancyhead[L]{\small\itshape \coursename}
\fancyhead[R]{\small\itshape \rightheader}
\fancyfoot[C]{\thepage}
\renewcommand{\headrulewidth}{0.4pt}
\renewcommand{\headrule}{\hbox to\headwidth{\color{rulecol}\leaders\hrule height \headrulewidth\hfill}}

% ----- Заголовки -----
\titleformat{\section}{\large\bfseries\color{accent}}{\thesection.}{0.6em}{}
\titleformat{\subsection}{\bfseries}{\thesubsection.}{0.5em}{}

% ----- Теоремные окружения (стиль конспекта) -----
\theoremstyle{definition}
\newtheorem{definition}{Определение}[section]
\newtheorem{example}[definition]{Пример}
\newtheorem{remark}[definition]{Замечание}
\theoremstyle{plain}
\newtheorem{theorem}[definition]{Теорема}
\newtheorem{lemma}[definition]{Лемма}
\newtheorem{corollary}[definition]{Следствие}
\newtheorem{property}[definition]{Свойство}
\newtheorem{criterion}[definition]{Критерий}
\renewcommand{\proofname}{Доказательство}

% ----- Операторы и сокращения -----
\DeclareMathOperator{\tg}{tg}
\DeclareMathOperator{\ctg}{ctg}
\DeclareMathOperator{\arctg}{arctg}
\DeclareMathOperator{\arcctg}{arcctg}
\DeclareMathOperator{\sh}{sh}
\DeclareMathOperator{\ch}{ch}
\DeclareMathOperator{\Th}{th}
\DeclareMathOperator{\cth}{cth}
\DeclareMathOperator{\Const}{const}
\DeclareMathOperator{\sign}{sign}
\DeclareMathOperator{\rang}{rang}
\DeclareMathOperator{\diam}{diam}
\DeclareMathOperator{\Span}{span}
\DeclareMathOperator{\Ker}{Ker}
\DeclareMathOperator{\Img}{Im}
\DeclareMathOperator{\tr}{tr}
\newcommand{\R}{\mathbb{R}}
\newcommand{\C}{\mathbb{C}}
\newcommand{\N}{\mathbb{N}}
\newcommand{\Z}{\mathbb{Z}}
\newcommand{\Q}{\mathbb{Q}}
\newcommand{\dx}{\,dx}
\newcommand{\dt}{\,dt}
\newcommand{\du}{\,du}
\newcommand{\eps}{\varepsilon}
\renewcommand{\Re}{\operatorname{Re}}
\renewcommand{\Im}{\operatorname{Im}}
\renewcommand{\le}{\leqslant}
\renewcommand{\ge}{\geqslant}
\newcommand{\scal}[2]{\left( #1,\,#2\right)}
\DeclarePairedDelimiter{\abs}{\lvert}{\rvert}
\DeclarePairedDelimiter{\norm}{\lVert}{\rVert}

\begin{document}
\begin{center}
  {\LARGE\bfseries \doctitle}\\[4pt]
  {\large\itshape \coursename}
\end{center}
\vspace{6pt}

\section{Билинейные и квадратичные формы}

\begin{definition}
\emph{Билинейная форма} — отображение $B\colon V\times V\to\R$, линейное по каждому аргументу. В
базисе $\{e_i\}$ её \emph{матрица} $A=(a_{ij})$, $a_{ij}=B(e_i,e_j)$, и
$B(x,y)=X^{T}AY$, где $X,Y$ — столбцы координат.
\end{definition}

\begin{theorem}[закон изменения матрицы]
При замене базиса с матрицей перехода $C$ ($X=CX'$) матрица формы преобразуется как
\[
  A'=C^{T}A\,C .
\]
\end{theorem}

\begin{proof}
$B(x,y)=X^{T}AY=(CX')^{T}A(CY')=X'^{T}(C^{T}AC)Y'$, поэтому матрица в новом базисе есть $C^TAC$.
(Это \emph{конгруэнтность}, в отличие от подобия $S^{-1}AS$ для операторов.)
\end{proof}

\begin{definition}
Форма \emph{симметрична}, если $B(x,y)=B(y,x)$ (т.е.\ $A=A^{T}$). \emph{Квадратичная форма}
$Q(x)=B(x,x)=X^{T}AX$. По симметричной $B$ форма $Q$ определяется однозначно, и обратно $B$
восстанавливается \emph{формулой поляризации}
\[
  B(x,y)=\tfrac12\bigl(Q(x+y)-Q(x)-Q(y)\bigr).
\]
\end{definition}

\section{Канонический вид. Закон инерции}

\begin{definition}
\emph{Канонический вид} квадратичной формы — диагональный: подходящей заменой координат
$Q=\sum_i \lambda_i z_i^2$. \emph{Нормальный} вид —
$Q=z_1^2+\dots+z_p^2-z_{p+1}^2-\dots-z_{p+q}^2$. Числа $p$ (положительный индекс), $q$
(отрицательный), $r=p+q$ (ранг) и $\sigma=p-q$ (сигнатура).
\end{definition}

\begin{theorem}[закон инерции Сильвестра]
Числа $p$ и $q$ не зависят от способа приведения к каноническому виду.
\end{theorem}

\begin{proof}[Инвариантность положительного индекса]
Покажем, что $p$ равно максимальной размерности подпространства, на котором $Q$ положительно
определена; это описание не ссылается на базис, поэтому $p$ инвариантно.

Пусть в каком-то нормальном базисе $L_+=\Span(e_1,\dots,e_p)$ — на нём $Q>0$ (для ненулевого
вектора), значит существует подпространство размерности $p$ с $Q>0$. Допустим, нашлось
подпространство $M$ большей размерности $\dim M>p$ с $Q>0$ на нём. Рассмотрим
$L_-=\Span(e_{p+1},\dots,e_n)$ (размерности $n-p$), где $Q\le0$. Тогда
$\dim M+\dim L_-=\dim M+(n-p)>n$, поэтому $M\cap L_-\ne\{0\}$: найдётся ненулевой вектор, на котором
одновременно $Q>0$ (как в $M$) и $Q\le0$ (как в $L_-$) — противоречие. Значит максимум равен $p$, и
$p$ — инвариант. Аналогично инвариантно $q$ (или из инвариантности ранга $r$).
\end{proof}

\section{Знакоопределённость квадратичных форм}

\begin{definition}
$Q$ \emph{положительно определена}, если $Q(x)>0$ при $x\ne0$ ($p=n$); \emph{отрицательно
определена}, если $Q(x)<0$ ($q=n$); \emph{знакопеременна}, если принимает значения обоих знаков
($p>0$ и $q>0$).
\end{definition}

\begin{theorem}[критерий Сильвестра]
Симметрическая форма с матрицей $A$ положительно определена тогда и только тогда, когда все
\emph{угловые} (главные) миноры положительны:
\[
  \Delta_1=a_{11}>0,\quad\Delta_2=\begin{vmatrix}a_{11}&a_{12}\\a_{21}&a_{22}\end{vmatrix}>0,\ \dots,\ \Delta_n=\det A>0 .
\]
Для отрицательной определённости знаки миноров чередуются: $(-1)^k\Delta_k>0$.
\end{theorem}

\begin{proof}[Идея (метод Якоби)]
Если все $\Delta_k\ne0$, метод Якоби приводит форму к каноническому виду
\[
  Q=\Delta_1 z_1^2+\frac{\Delta_2}{\Delta_1}z_2^2+\dots+\frac{\Delta_n}{\Delta_{n-1}}z_n^2 ,
\]
где канонические коэффициенты суть отношения соседних угловых миноров. Положительная определённость
равносильна положительности всех коэффициентов, т.е.\ $\tfrac{\Delta_k}{\Delta_{k-1}}>0$ при всех
$k$ ($\Delta_0=1$), что по индукции эквивалентно $\Delta_k>0$ для всех $k$. Для отрицательной
определённости все коэффициенты отрицательны, что даёт чередование знаков $\Delta_k$.
\end{proof}

\end{document}
